% packing-and-marker-examples.tex — 打包与 Marker 算例
% Source: examples/packing-and-marker-examples.md

\section{打包与 Marker 算例}\label{sec:packing-examples}

对应章节：\cref{ch:packing}

\subsection{算例 1：单组系数的 data packing}

\begin{examplebox}[title={Packing Worked Example：单组系数的 data packing}]

\textbf{输入}：1 个 band 行，$N_g = 4$，$F_s = 0$（jointly），$\mathit{GTLI} = 2$。

4 个 sign-magnitude 系数（已左移回原始 bitplane 位置）：

\begin{tabular}{@{}rrrrl@{}}
\toprule
索引 & 系数值 & 符号 & 幅度 & 二进制 \\
\midrule
0 & 0x0000001C & $+$ & 28 & 0b00011100 \\
1 & 0x00000005 & $+$ & 5  & 0b00000101 \\
2 & 0x8000000A & $-$ & 10 & 0b00001010 \\
3 & 0x00000000 & $+$ & 0  & 0b00000000 \\
\bottomrule
\end{tabular}

$\mathit{GCLI} = \mathrm{BSR}(28\,|\,5\,|\,10\,|\,0) + 1 = \mathrm{BSR}(31) + 1 = 4 + 1 = 5$。

\texttt{pack\_data()} 执行：

\textbf{Step 1：写符号位}（$F_s=0$，每个 bitplane 前写 4 个符号 bit）：
\begin{lstlisting}[numbers=none]
sign bits: 0, 0, 1, 0  （系数 0 正, 1 正, 2 负, 3 正）
\end{lstlisting}

\textbf{Step 2：bitplane 4（MSB）}——从 $\mathit{gcli}-1 = 4$ 开始：
\begin{lstlisting}[numbers=none]
bp=4: buf[0]>>4 = 0, buf[1]>>4 = 0, buf[2]>>4 = 0, buf[3]>>4 = 0
写入: 0 0 0 0
\end{lstlisting}

\textbf{Step 3：bitplane 3}：
\begin{lstlisting}[numbers=none]
bp=3: buf[0]>>3 = 1, buf[1]>>3 = 0, buf[2]>>3 = 1, buf[3]>>3 = 0
写入: 1 0 1 0
\end{lstlisting}

\textbf{Step 4：bitplane 2}（$\mathit{GTLI} = 2$，这是最后一个写入的 bitplane）：
\begin{lstlisting}[numbers=none]
bp=2: buf[0]>>2 = 1, buf[1]>>2 = 1, buf[2]>>2 = 0, buf[3]>>2 = 0
写入: 1 1 0 0
\end{lstlisting}

\textbf{完整 bit 流}（1 组）：
\begin{lstlisting}[numbers=none]
符号: 0 0 1 0  |  bp=4: 0 0 0 0  |  bp=3: 1 0 1 0  |  bp=2: 1 1 0 0
\end{lstlisting}

总计：$4 \text{ (sign)} + 4 \times 3 \text{ (data)} = \mathbf{16 \text{ bit}} = 2 \text{ bytes}$。

验证：$(\mathit{GCLI} - \mathit{GTLI} + 1) \times \mathit{group\_size} = (5 - 2 + 1) \times 4 = 16 \text{ bit}$ \checkmark

\end{examplebox}
